\subsubsection{LGV 引理 (Weighted LGV Lemma)}

LGV 引理，适用于 \textbf{DAG} 上带权不相交路径问题。
设边权为 $w(e)$（可为数值或多项式），定义路径 $P$ 的权值为 $\omega(P) = \prod_{e \in P} w(e)$。
构造矩阵 $M$，其中 $M_{ij}$ 为起点 $a_i$ 到终点 $b_j$ 的所有路径权值之和：

$$
M_{ij} = \sum_{P: a_i \to b_j} \omega(P) = \sum_{P: a_i \to b_j} \prod_{e \in P} w(e)
$$

若图满足条件（如平面图且起点终点排序使得仅 $\sigma=id$ 存在不相交路径），则所有不相交路径组 $(P_1, \dots, P_n)$ 的权值积之和为：

$$
\det(M) = \sum_{(P_1, \dots, P_n) \in \text{Disjoint}} \prod_{k=1}^n \omega(P_k)
$$

\noindent \textbf{注}：若求不相交路径总数，设 $w(e)=1$；若求路径权值总和，设 $w(e)=x^{val}$ 并做多项式运算。

\paragraph{使用条件与多测} 一般 LGV 行列式给出带排列符号的点不相交路径族和；只有证明非恒等配对均不可能时才等于无符号计数，平面性本身不充分。这里不相交指顶点不相交。每组重建 DAG 路径 DP 和 M；对各源沿拓扑序传播，别把每条路径显式枚举。环图不能直接按本式做有限路径和。
